\section{Two exact improvements to executable packing kernels}\label{sec:kernels}
The schedule theorems concern asymptotic contraction certificates. The next
results operate on explicit finite graphs and do not need those asymptotic
packing margins to hold. They preserve the corresponding original routines'
choices exactly, so any valid prior packing certificate continues to apply.
Incremental conditional expectations and common-denominator arithmetic are
classical techniques; our novelty claim is the specified implementation and
proved improvement relative to the reviewed executable corpus.

\subsection{Incremental products for matching potentials}
Partition the bipartite edge occurrences into matchings $E_1,\ldots,E_K$.
Parallel edges remain distinct and occur in different matchings. Each vertex
has one of $Q$ labels. For a partially assigned vertex $v$, let $p_v(a)$ be
its assigned-label indicator, or $1/Q$ if it remains unassigned. In this
subsection $\eta\ge1$ denotes the reference's concentration parameter, called
\texttt{lam} in code; it need not equal a tensor-size logarithm. Put
$t=1+1/\eta$ and let $h_j$ be the reference thresholds. Define
\begin{equation}\label{eq:matching-products}
 P_{j,a,b}=t^{-h_j}\prod_{(u,v)\in E_j}
 \left(1+\frac{p_u(a)p_v(b)}{\eta}\right),\qquad
 \Phi=\sum_{j,a,b}P_{j,a,b}.
\end{equation}
Because edges within a matching have disjoint endpoints, this is exactly the
conditional expectation of $\sum_{j,a,b}t^{X_{j,a,b}-h_j}$. Independence
between different matchings is unnecessary: expectations add. The latest
matching packet evaluates~\eqref{eq:matching-products} afresh for every
candidate label at every vertex~\cite{matching}.

\begin{theorem}[Same choices with incremental matching products]\label{thm:matching}
For an explicit graph with $M$ edge occurrences and $V$ incident vertices,
the implementation \texttt{label\_graph\_incremental} returns the same
first-fit matching partition, labels, thresholds, loads, and complete selected
potential trace as the pinned reference. Its update phase uses
$O(MQ^2+VQ)$ rational operations, compared with $\Theta(VMQ^3)$ elementary
product operations in the nonempty reference's repeated evaluations.
Initialization, coloring, bit costs, and extra storage are charged below.
\end{theorem}
\begin{proof}
Suppose an unassigned left vertex $u$ is processed. Each incident matching
contains exactly one edge $(u,v)$. For every opposite label $b$ with
$p_v(b)>0$, set
\[
 f_b=1+\frac{p_v(b)}{Q\eta},\qquad R_a=P_{j,a,b}/f_b.
\]
The divisor is positive. In this slice the old contribution is
$\sum_aP_{j,a,b}$. If $u$ receives candidate label $c$, the new contribution is
\[
 \sum_aR_a+\frac{p_v(b)}\eta R_c.
\]
Thus all candidates share a base, and only the last correction depends on
$c$. Summing affected slices gives
\begin{align}
 B&=\Phi+\sum_{\text{affected slices}}
          \left(\sum_aR_a-\sum_aP_{j,a,b}\right),\notag\\
 S_c&=\sum_{\text{affected slices}}\frac{p_v(b)}\eta R_c,
 &\Phi_c&=B+S_c.\label{eq:matching-update}
\end{align}
No slice is changed twice within a matching because $E_j$ is a matching.
The right-vertex rule interchanges the indices. Dividing the affected table
entries by their old factors, then multiplying the chosen entries by
$1+p_v(b)/\eta$, restores~\eqref{eq:matching-products} exactly. The averaging
identity $\sum_c\Phi_c=Q\Phi$ follows either from the formulas or from
conditional expectation and is checked exactly by the implementation.

Candidate values and smallest-label tie breaking therefore equal those from
recomputation. Induction gives every selected trace value and final load.
The endpoint-occupancy sets used in first-fit coloring make the same first
available matching choice as the reference's edge scans; input order is
retained. This proves full-output equality, including tuple-pair and list-pair
edge representations.

When an edge's first endpoint is processed, its other endpoint is uniform;
there are $Q$ affected slices of $Q$ cells. When the second endpoint is
processed, there is one affected slice of $Q$ cells. Summing over edge
occurrences gives $O(M(Q^2+Q))$ rational operations, plus $O(VQ)$ candidate
summation and selection operations. In contrast, the reference has $VQ$
full potential evaluations, each scanning $M$ edge factors in $Q^2$ cells.
\end{proof}

\paragraph{Costs outside the update phase.}
The initial value of every cell for matching $j$ is
\[
 t^{-h_j}\left(1+\frac1{\eta Q^2}\right)^{|E_j|}.
\]
Binary powering costs
$O(K(\log(M+1)+\log(h_{\max}+1)))$ rational multiplications; the tables
require $O(KQ^2)$ initial assignments. First-fit coloring uses $O(MK)$
endpoint-set membership tests, and vertex sorting uses $O(V\log V)$
comparisons. The final product-invariant check costs $O(KQ^2)$.
The maintained tables contain $KQ^2$ rational values, in addition to graph
information, $Q$ scores, and $V$ selected trace values. This deliberately
spends more product storage than recomputation, while remaining polynomial
in the application parameters.

A common denominator for the conditional products divides
$(\eta+1)^{h_{\max}}\eta^{m_{\max}}Q^{2m_{\max}}$, where
$m_{\max}=\max_j|E_j|$. Product bit lengths are therefore
$O((h_{\max}+m_{\max})\log(\eta Q+2))$, with logarithmic overhead for sums.
Multiplication, division, and greatest-common-divisor costs must be paid at
those lengths. The theorem's operation count does not assume constant-time
large-integer arithmetic. Bit sizes are polynomial for explicit graphs and
the concentration parameters used by the reviewed contraction schedules.

\subsection{Common-denominator integer scores for column balancing}
The two-stage exponential-potential construction first groups rows, producing
nonnegative integer weights $w_{r,v}$ from row group $r$ to column vertex $v$.
There are $Q$ row groups, $Q$ possible column labels, and $N$ column vertices.
Let $z=a/b>1$ be in lowest terms. The earlier packet refinement already caches
suffix products and row-power sums, achieving $O(NQ^2)$ dense arithmetic
work~\cite{packet}. We do not claim that order or the idea of caching as new.
Our change removes repeated rational normalization from candidate scoring.

At time $t$, let $L_{r,c}$ be the assigned weight in cell $(r,c)$,
$S_r=\sum_cL_{r,c}$, $W_r=\sum_vw_{r,v}$, and $W=\max_rW_r$. Set
\begin{align}
 C_w&=(Q-1)b^w+a^w,\notag\\
 A_r&=b^{W-W_r}\prod_{v\text{ unassigned}}C_{w_{r,v}},
 &H_{r,c}&=a^{L_{r,c}}b^{S_r-L_{r,c}}.\label{eq:integer-state}
\end{align}
All these quantities are nonnegative integers. The unnormalized conditional
potential has the exact representation
\begin{equation}\label{eq:integer-invariant}
 \Phi=\frac{\sum_{r,c}A_rH_{r,c}}{Q^{N-t}b^W}.
\end{equation}
Indeed the ordinary row-cell expectation is
$z^{L_{r,c}}\prod_{v\text{ unassigned}}(Q-1+z^{w_{r,v}})/Q$;
clearing the common denominator gives~\eqref{eq:integer-invariant}.
An optional common threshold multiplier $z^{-B_0}$ has no effect on choices.

For the next column, put $w_r=w_{r,v}$ and divide
$A'_r=A_r/C_{w_r}$ exactly. Candidate $c$ then has numerator
\begin{equation}\label{eq:integer-candidate}
 D_c=\sum_r A'_r\left[
 b^{w_r}\sum_hH_{r,h}+(a^{w_r}-b^{w_r})H_{r,c}\right]
\end{equation}
over the common denominator $Q^{N-t-1}b^W$. Moreover
\[
 \sum_cD_c=\sum_{r,c}A_rH_{r,c}.
\]
Choose the least integer numerator, breaking ties by label index. Multiply
the chosen $H_{r,c}$ by $a^{w_r}$ and each other $H_{r,h}$ by $b^{w_r}$.
These updates restore~\eqref{eq:integer-state}. By induction, every chosen
label and every reconstructed rational candidate potential equals the
reference. Zero weights, empty columns, and $Q=1$ obey the same identities.

Dense arithmetic remains $O(NQ^2)$. The core score computations use integers
and divisions known to be exact; rational reduction is deferred to the
requested diagnostic output. The supplied implementation still constructs
the full candidate trace, so its measured gains do not arise from omitting
that work. It removes the suffix-product table and uses $O(Q^2)$ main update
state, excluding input, power dictionaries, and the explicitly requested
$O(NQ)$ candidate output objects.

Since $C_w\le Qa^w$, the main numerator bit length is at most
$O(N\log Q+W\log a+\log Q)$. Threshold normalization adds
$O(B_0\log a)$ bits. Exact integer divisions and output rational reductions
are charged normally. These bounds concern explicit edge multiplicities;
we make no polynomial-in-log-weight claim for compressed enormous weights.
Unlike the matching theorem, this result does not assert a better dense
asymptotic operation count than the strongest cached reference. It supplies
an exact representation change and a measured implementation gain.
